%! The Inverse and Determinant of a Matrix — Unit 4, Lesson 4.11 — Exit Ticket (Answer Key)
\documentclass[10pt]{article}
\usepackage{precalculus-article}
\usepackage{precalculus-key}   % pulls in -boxes; do NOT also load -boxes
\newcommand{\TallMath}[1]{$\displaystyle #1\rule[-1.4em]{0pt}{3.2em}$}

\begin{document}

\pageheader{Unit 4, Lesson 4.11}{Exit Ticket --- Answer Key}
\namedateperiod

\vspace{0.12in}

\begin{notesbox}{Exit Ticket}
\begin{enumerate}[leftmargin=1.5em, itemsep=6pt]

    \item Compute $\det\begin{bmatrix}5&2\\3&1\end{bmatrix}$.\\[4pt]
          \ans{$\det = 5(1) - 2(3) = 5 - 6 = -1$}
          \vspace{0.25in}

    \item Does the matrix in question 1 have an inverse? If so, find it.\\[4pt]
          \ans{Yes, since $\det = -1 \ne 0$.}\\[4pt]
          \ans{$A^{-1} = \dfrac{1}{-1}\begin{bmatrix}1&-2\\-3&5\end{bmatrix}
          = \begin{bmatrix}-1&2\\3&-5\end{bmatrix}$}
          \vspace{0.25in}

    \item A $2\times2$ matrix has $\det = 0$. What does this mean geometrically
          about its two column vectors?\\[4pt]
          \ansline{The two column vectors are parallel (one is a scalar multiple of the other).}
          \ansline{They span only a line, not a plane; the parallelogram they form has area zero.}

\end{enumerate}
\end{notesbox}

\end{document}
